% Abstrat of your communication to the
% ACA2019 Conference (ÉTS, Montréal, Canada, 2019).
% You must use LaTeX format.


%====================================================
% IMPORTANT: DO NOT MODIFY THE FOLLOWING LINES
%====================================================

\documentclass{article}

\usepackage[T1]{fontenc}
\usepackage[utf8]{inputenc}
\usepackage{amssymb,amsthm,amsmath, amsfonts, latexsym}
\usepackage{graphicx}
\usepackage{hyperref}
\usepackage{times}

\setlength{\oddsidemargin}{1.46cm}
\setlength{\textwidth}{13cm}
\setlength{\topmargin}{0.0cm}
\setlength{\textheight}{20.5cm}

\makeatletter
\newcommand{\TITLE}[2][]{
  \begin{center} \Large \bfseries #2%
  \def\@temp{#1}\ifx\@temp\@empty\relax\else\footnote{#1}\fi\end{center}}

\newcommand{\AUTHOR}[2][]{\textbf{#2}%
  \def\@temp{#1}\ifx\@temp\@empty\relax\else\textsuperscript{#1}\fi}

\newcommand{\ADDRESS}[2]{\noindent
  \textsuperscript{#1}\parbox[t]{0.9\textwidth}{#2} \vspace{6pt}}

\makeatother

\newcommand{\KEYWORDS}[1]{\paragraph{Keywords:} #1}
%\newcommand{\MSCLASSIFICATION}[1]{\paragraph{Mathematics Subject Classification 2010:} #1}

\renewcommand{\thefootnote}{\fnsymbol{footnote}}
\setcounter{footnote}{0}
\pagestyle{empty}

\newcommand{\HEADING}{
{\noindent \small Applications of Computer Algebra -- ACA2019 \\
\'Ecole de technologie sup\'erieure, July 16--20, 2019} \vspace{10pt}}

%===============================================
% DO NOT MODIFY THE ABOVE LINES
%===============================================

\begin{document}

\HEADING


\TITLE{Computer tools as tools of semiotic mediation in studying infinitesimal analysis: The didactical challenge}

% If you need to thank someone next to the title, please use the next line (and delete the previous one)

% \TITLE[thanks]{Title of talk}
% Write the author names in the order you wish them to appear
% and underline the one who will give the talk.
% Identify each one of them with a number
% in order to give their own affiliations later on. 

\begin{center}
 \AUTHOR[1]{\underline{Anatoli Kouropatov}} % the author who will give the talk
% IF two or more authors have the same affiliation, associate to them the same number.
% For instance, if there are three authors and the first a the second ones have the same affiliation,
% write as follows:
% \AUTHOR[1]{Third Author}
\end{center}


% Write now the abstract of your communication. 

The use of computer technology can offer a lot to math education. The
interactivity and dynamics of technology can enable the creation of
representations and tools that allow teachers to present and demonstrate
mathematical ideas and concepts. These tools also allow students to
explore these concepts and reach conclusions about their properties and
relations between them by considering the effect of the manipulations
performed on them in a certain virtual environment (Moyer-Packenham \&
Westenskow, 2013).

I argue that these digital representations and tools become cognitive
artifacts (in the Vigotskian perspective) used by learners as tools
of semiotic mediation (Bartolini Bussi \& Mariotti, 2008). These may be
determined by the specific interface of a certain digital environment,
as well as by the learner’s personal perception.

From my research (which was conducted with first year engineering
students) I observed that as a result of that mediation, students
(sometimes) construct a vague and inconsistent world of mathematical
meanings. I speculate that to overcome this problem we need an appropriate
didactical design of computer-based tasks that allow (through the use
of specific artifacts in the mediation process) the construction of the
appropriate mathematical meaning in students' minds.

In the conference I will demonstrate the empirical evidence of students'
work in a digital environment that leads them to inconsistent and
problematic mathematics concepts, and I will discuss how to identify
the origin of students' difficulties as well as how to design digital
tasks to overcame (or even prevent) these difficulties in the case of
studying infinitesimal analysis.
 

% Write the keywords separated by commas

\KEYWORDS {Computer tools, semiotic mediation, infinitesimal analysis, task design }

% Write the MSC2010 codes corresponding to your communication

%\MSCLASSIFICATION{Code1, Code2}



% Next include the references you need. 
% If your abstract includes no references at all, please delete all the lines
% from \begin{thebibliography} to \end{thebibliography}.

% Please, do not use more than 9 references.

\begin{thebibliography}{9}

\bibitem{1} \textsc{Moyer-Packenham, P.S.; Westenskow, A.},
\emph{Effects of virtual manipulatives on student achievement and mathematics
learning}.
\emph{International Journal of Virtual and Personal Learning Environments
(IJVPLE)}, \textbf{4}(3), 35--50, (2013).

\bibitem{2} \textsc{Bartolini Bussi, M. G.; Mariotti, M. A.},
\emph{Semiotic mediation in the mathematics classroom: Artifacts and signs
after a Vygotskian perspective}.
In \emph{Handbook of International Research in Mathematics Education (2nd ed.)},
L. English; M. Bartolini Bussi; G. Jones; R. Lesh; D. Tirosh (eds.),
746–-783, Mahwah, NJ: Lawrence Erlbaum, 2008.

  % Model for references to books.
%  \bibitem{label1} \textsc{A. Uthor1; A. Uthor2}, \emph{Title}.
%    Editorial, City, year of publication.

 %  Model for references to articles
%  \bibitem{label2} \textsc{A. Uthor}, Title.
%  \emph{Journal} \textbf{volume}(number), first page--last page  (year of publication).

  %  Model for references to proceedings. 
%  \bibitem{label3} \textsc{A. Uthor}, Title.
%  In \emph{Title of the proccedings}, names of the editors (eds.),
%   first page--last page. Editor(s), City, year of publication.

\end{thebibliography}

\vspace{1cm}

% Full Affiliation of all authors
\ADDRESS{1}{
 Mathematics Department\\Levinsky College of Education \\
Tel-Aviv, Israel \\
\texttt{anatoliko@gmail.com}
}

%\vspace{.5cm}
%
%\ADDRESS{2}{
%Department1  \\University1 \\
%Posta Address1 \\
%\texttt{e@mail1}
%}


\end{document}
